% !TeX root = ../nc-del-pezzo.tex
% 2026-11-20 JST / v0.1
\subsection{具体的量子化による全 Artin 変形の表示}
\label{subsec:quantization-formal-universal}

滑らかな可換点 $X$ を固定し，グラフ formality 写像を一つ選ぶ．
定理\ref{thm:quantization-descent}の構成を特殊ファイバーの同一視を保つ
変形に制限すると，全ての局所 Artin $\mathbb C$-代数の射に自然な写像
\[
 Q_X:\mathrm{PoisDef}_X\longrightarrow\mathrm{Def}^{\mathcal O}_X
\]
を得る．これは stack 化後だけの写像ではない．Artin 局所基底上では
基本線束の持上げと重み付き生成元を大域的に選べ，その選択を変える
正規形の座標変換は実際の次数付き代数同型を与える．
次数スカラーは構造対象を保つ自然同型で恒等化されるから，
$Q_X$ は定義\ref{def:framed-o-abel}の同型類に値を取る．
ここで標準形の選択を与える滑らかな座標変換群の torsor は，
特殊ファイバーの $\mathbb C$-点を冪零厚化へ持ち上げることで自明化できる．

\begin{lemma}[完備局所環の自己写像]
\label{lem:formal-endomorphism-inverse}
$(C,\mathfrak m)$ を剰余体 $k$ を持つ完備 Noether 局所 $k$-代数とする．
連続局所 $k$-代数自己準同型 $\varphi:C\to C$ が
$\mathfrak m/\mathfrak m^2$ 上で同型を誘導するなら，$\varphi$ は同型である．
\end{lemma}
\begin{proof}
$\operatorname{gr}_{\mathfrak m}C$ は次数1部分で生成されるため，
$\varphi$ は各 $\mathfrak m^r/\mathfrak m^{r+1}$ 上で全射を誘導する．
これらは有限次元 $k$-ベクトル空間だからその写像は同型である．
帰納的に $C/\mathfrak m^r$ 上の写像も同型となり，逆極限を取ればよい．
この議論は $C$ の正則性を仮定しない．
\end{proof}

\begin{proposition}[具体的量子化の接写像]
\label{prop:quantization-tangent-isomorphism}
$V_X=H^0(X,\bigwedge^2T_X)$ と置く．一次 commutator は完全列
\[
 0\longrightarrow H^1(X,T_X)
 \longrightarrow T\mathrm{Def}^{\mathcal O}_X
 \xrightarrow{\,\beta\,}V_X\longrightarrow0
\]
を与える．$TQ_X$ は同型であり，左辺の可換方向と右辺の初括弧の双方で
恒等写像を誘導する．
\end{proposition}
\begin{proof}
$R_\epsilon=\mathbb C[\epsilon]/(\epsilon^2)$ とする．
命題\ref{prop:abel-poisson-formal-comparison}の証明で用いた
Lowen の比較と $H^1(\mathcal O_X)=H^2(\mathcal O_X)=0$ による untwisting により，
Abel 一次変形の同型類は普通の結合代数層変形の gauge 同値類に一致する
\cite[Theorems~3.14,3.26]{Lowen2008}
\cite[Corollary~7.5]{Yekutieli2015}．
その代表 $\mathcal A_\epsilon$ に対して
\[
 \epsilon\,\beta(a,b)
   =[\widetilde a,\widetilde b]_{\mathcal A_\epsilon}
\]
と定める．平坦性による $\epsilon\mathcal A_\epsilon\simeq\mathcal O_X$ を用いており，
局所的な持上げの選択によらない．結合性から $\beta$ は二重導分で，
gauge 同値によって変化しないため，$V_X$ の元を与える．
この構成は任意の平方零係数空間に関して自然であり，接空間上で線型である．
$\beta=0$ は $\mathcal A_\epsilon$ が可換であることと同値なので，
その核は通常のスキーム一次変形 $H^1(X,T_X)$ である．
可換代数層間の結合代数同型は可換代数同型でもあるから，左の写像は単射である．

一方，零 Poisson 構造における接空間は
$T\mathrm{PoisDef}_X=H^1(X,T_X)\oplus V_X$ である．
純粋な $\epsilon\pi$ に対する具体的量子化は
\[
 a\star b=ab+\frac\epsilon2\widetilde\pi(\mathrm da,\mathrm db),
 \qquad\Omega=F.
\]
$F$ は Casimir なので $F\star a=Fa$ であり，中心商の一次積は
$B_0\oplus\epsilon B_0$ 上の同じ式である．
$B_0$ の非零斉次元の持上げ $f$ を invert する際，
平坦性と $B_0$ の整域性から $f$ は左右正則である．
また二重 commutator が零なので
$b f^{k+1}=f^k(fb+(k+1)[b,f])$ であり，
$f$ の冪は左右の Ore 条件を満たす．その局所化の一次積は，
二重導分を通常の商則で延長した式になる．従って標準アフィン開
$D_+(f)$ の重み零部分では初括弧は $\pi|_{D_+(f)}$ に等しい．
次数2では $D_+(x_i)$，次数1ではこれらと $D_+(yz)$ が $X$ を被覆する．
前者の $x_i$，後者の $z\star y^{-1}$ は局所化後の次数1の可逆元であり，
この可逆性は冪零変形でも保たれる．従って局所化は強次数付きとなり，
その次数零加群圏が tails の対応する局所化に一致する．
よって $\beta(TQ_X(0,\pi))=\pi$ である．
とくに $\beta$ は全射である．

Poisson 構造が零なら積も中心元も可換のものに一致し，
$Q_X$ は通常の $X_R$ を文字通り
$(\operatorname{Qcoh}(X_R),\mathcal O_{X_R})$ に送る．
よって $TQ_X$ は完全列の核と商の双方で恒等写像を誘導する．
二つの短完全列の比較から $TQ_X$ は同型である．
\end{proof}

\begin{remark}[Hochschild 接空間との一致]
一次 Abel 変形の接空間は Hochschild コホモロジーの二次群であり
\cite[Theorem~3.1]{LowenVdBHochschild}，
準連接層の Hochschild コホモロジーはスキームのものに一致する
\cite[Corollaries~7.7.2,7.8.2]{LowenVdBHochschild}．
HKR 分解\cite[Corollary~0.6]{YekutieliHKR2002}と $H^2(X,\mathcal O_X)=0$ により
\[
 T\mathrm{Def}^{\mathcal O}_X
 \simeq\operatorname{HH}^2(X)
 \simeq H^1(X,T_X)\oplus H^0(X,\bigwedge^2T_X).
\]
上の証明は，この分解の積構造や Gerstenhaber 括弧との両立を必要としない．
初 commutator によって右側を正規化すれば，$TQ_X$ は対角成分が恒等の
上三角行列となる．
\end{remark}

\begin{theorem}[具体的量子化の全 Artin 普遍性]
\label{thm:concrete-quantization-all-artin}
固定したグラフ formality 写像に対する $Q_X$ は自然な全単射である．
従って全ての framed な $\mathcal O$-marked 平坦 Abel Artin 変形は，
本節の重み付き多項式環の量子化の中心超曲面商の tails として実現できる．
\end{theorem}
\begin{proof}
命題\ref{prop:abel-poisson-formal-comparison}の既知の自然な全単射の逆を
$\chi:\mathrm{Def}^{\mathcal O}_X\xrightarrow{\sim}\mathrm{PoisDef}_X$ とする．
定理\ref{thm:abel-formal-ring}の証明により後者は表示された完備 Noether 局所環
$R_X$ で prorepresent される．従って自然変換 $\chi Q_X$ は
連続局所 $\mathbb C$-代数自己準同型 $\varphi:R_X\to R_X$ に対応する．
命題\ref{prop:quantization-tangent-isomorphism}と $T\chi$ の可逆性から
$\varphi$ の余接写像は同型である．
補題\ref{lem:formal-endomorphism-inverse}により $\varphi$ は同型となり，
$Q_X$ は全ての Artin 環上で全単射である．
$Q_X$ の構成は，まさに主張した中心商の tails を与える．
\end{proof}

\begin{corollary}[通常の次数付き中心超曲面表示]
\label{cor:all-artin-weighted-hypersurface-presentation}
任意の局所 Artin $R$ と任意の
$[\mathcal D,\mathcal O_R]\in\mathrm{Def}^{\mathcal O}_X(R)$ に対して，
通常の非負次数付き $R$-代数 $A_R$ と正則中心元 $\Omega\in(A_R)_m$ が存在し，
\[
 (\mathcal D,\mathcal O_R)
 \simeq(\operatorname{QGr}(A_R/(\Omega)),\pi(A_R/(\Omega)))
\]
である．$A_R$ は可換点の重み付き多項式環の平坦変形であり，
命題\ref{prop:nilpotent-weighted-as}と
命題\ref{prop:flat-central-hypersurface-ext}の相対的結論を満たす．
特に，中心商 $B_R=A_R/(\Omega)$ について左右とも
$\underline{\operatorname{Ext}}_{B_R}^{n+1}(R,B_R)=R(n-1)$ で，
他の augmentation Ext は零である．
\end{corollary}
\begin{proof}
定理\ref{thm:concrete-quantization-all-artin}によって量子化表示を取り，
定理\ref{thm:nilpotent-universal-quantization}の次数別自由性と
普遍的中心正則性を用いる．これらは上記二命題の仮定を満たし，
$N=n+2$，$W-m=n-1$ を代入すると主張の Ext 公式を得る．
\end{proof}

この結論は中心超曲面表示の存在をいうもので，その表示の一意性をいうものではない．
$Q_X$ が既知の比較 $\chi^{-1}$ と一致することも必要ない．
full groupoidの同値には第\ref{sec:formal-groupoid}節の微小自己同値の剛性を，
相対Serre性には第\ref{sec:artin-ample-perfect}節と第\ref{sec:relative-serre-artin}節の
compact性・対角の持上げを用いる．大域ncモジュライの代数性と全既約成分数は未決定である．
